\documentclass[11pt]{amsart}
\usepackage{amssymb,mathtools}
\usepackage[margin=1.15in]{geometry}
\usepackage[hidelinks]{hyperref}
\newtheorem{theorem}{Theorem}[section]
\newtheorem{lemma}[theorem]{Lemma}
\theoremstyle{definition}\newtheorem{definition}[theorem]{Definition}
\theoremstyle{remark}\newtheorem{remark}[theorem]{Remark}
\newcommand{\UU}{\mathcal U}
\newcommand{\Hom}{\operatorname{Hom}}
\newcommand{\Nat}{\operatorname{Nat}}
\newcommand{\Spec}{\operatorname{Spec}}
\newcommand{\Zar}{\operatorname{Zar}}
\newcommand{\tr}[1]{\left\Vert #1\right\Vert}
\title[The essential image of the functor of points]{The essential image of the functor of points\\Compact opens and finite affine atlases}
\date{October 29, 2026}
\begin{document}\raggedbottom\maketitle

\section{The statement}
We use all commutative rings with set carriers in $\UU_\ell$ as affine
tests. Their object type lies in $\UU_{\ell+1}$ and their hom sets lie
in $\UU_\ell$. A functor takes values in sets in $\UU_\ell$.
All rings and ring maps are unital; the zero ring is allowed.
For a ring $A$ put $h_A(B)=\Hom(A,B)$.

We seek an intrinsic description of the essential image of
$X\mapsto h_X$, where $h_X(B)=\Hom(\Spec B,X)$ and $X$ is a small qcqs
locally ringed lattice. Full faithfulness is the preceding comparison
result. The new assertion is existence of a representing scheme.
We follow the compact-open formulation of \cite{Zeuner,ZH}, while
constructing the resulting lattice directly from finite chart data.

\begin{definition}
A functor $F$ is a \emph{Zariski sheaf} if compatible elements on every
finite basic cover of a test ring have a unique amalgam. Explicitly, if
$\bigvee_i D(a_i)=\top$ in $\Zar(B)$, then
\[
 F(B)\longrightarrow\prod_i F(B[1/a_i])
 \rightrightarrows\prod_{i,j}F(B[1/(a_i a_j)])
\]
is an equaliser, with the canonical maps of localisations. Empty covers
are included; in particular $F(0)$ is contractible.
\end{definition}
Different presentations of the localisations are related by their
universal properties. The sheaf condition can equivalently be stated
for arbitrary choices of the corresponding localisation rings.

The Zariski lattices define a functor $\Zar$ on rings. For $h:A\to B$,
its lattice map is characterised by
\[
 \Zar(h)(D(a))=D(h(a)).
\]
Uniqueness of this characterisation proves identity and composition.

\begin{definition}
A \emph{compact open} of $F$ is a natural transformation $U:F\to\Zar$.
Its realisation is the subfunctor
\[
 |U|(B)=\sum_{x:F(B)}(U_B(x)=\top).
\]
It is \emph{affine} if there merely exist a small ring $A$ and a natural
equivalence $h_A\simeq |U|$. A finite family $U_i$ covers $F$ if
$\bigvee_i U_i=\top$ pointwise. A \emph{functorial qcqs scheme} is a
Zariski sheaf merely admitting such a finite affine compact-open cover.
\end{definition}
Compactness here is a condition on every affine pullback. Thus the
compactness of the overlaps is built into the definition. We do not
require the overlaps to be affine, nor add that as a gluing hypothesis.

\begin{theorem}[Essential image]\label{thm:essential}
A small-set-valued functor $F$ is a functorial qcqs scheme if and only if
there merely exist a small qcqs locally ringed lattice $X$ and a natural
equivalence $F\simeq h_X$.
\end{theorem}
Together with full faithfulness, this gives the comparison on objects
and arrows. To state an equivalence of univalent categories one must
also package the geometric category and prove its univalence. That
packaging is separate from Theorem~\ref{thm:essential}.

\section{Compact opens on affine tests}
For any $F$, evaluation at $1_A$ gives the covariant Yoneda equivalence
\[
 \Nat(h_A,F)\simeq F(A).
\]
Its inverse sends $x$ to $h\mapsto F(h)(x)$. The two inverse laws are
the identity law and the naturality equation at $h:A\to B$.
In particular compact opens on $h_A$ correspond to elements of $\Zar(A)$.
For $u\in\Zar(A)$ the corresponding subfunctor is
\[
 B\longmapsto\{h:A\to B\mid\Zar(h)(u)=\top\}.
\]
If $u=D(a)$, its realisation is $h_{A[1/a]}$, by the universal property
of localisation and $D(h(a))=\top$ if and only if $h(a)$ is a unit.

More generally, for $x\in F(A)$, pullback of $|U|\to F$ along the
transformation $h_A\to F$ determined by $x$ is the affine compact open
classified by $U_A(x)$. This is a direct consequence of naturality:
\[
 U_B(F(h)(x))=\Zar(h)(U_A(x)).
\]
This formula is the precise link between the classifier definition and
compact open subfunctors. No ideal representatives are chosen.

The compact opens of $F$ form a distributive lattice with pointwise
operations. This lattice initially lies in $\UU_{\ell+1}$, because its
elements are natural transformations on all small rings. Even though
each value $\Zar(A)$ is small, the entire lattice must not simply be
asserted small.

\section{The geometric functor has an affine compact-open atlas}
Let $X$ be a small qcqs locally ringed lattice. For an open $e$ of $X$
define
\[
 U_e(p)=p^*(e),\qquad p:\Spec B\to X.
\]
This is natural: the inverse-image map of $\Spec h$ is $\Zar(h)$,
and inverse images compose. The slice universal property gives
\[
 |U_e|(B)\simeq\Hom(\Spec B,X|e).
\]
Indeed a map factors uniquely through $X|e$ exactly when $p^*(e)=\top$;
the factorisation preserves locality. Thus affine chart opens give
affine compact opens, and a finite chart cover remains a cover.

The functor $h_X$ is a Zariski sheaf: maps from the basic open pieces
of $\Spec B$ glue uniquely by the finite-cover descent theorem for
morphisms, using $(\Spec B)|D(a)\simeq\Spec(B[1/a])$. This includes
the empty-cover case. Consequently $h_X$ satisfies the left side of
Theorem~\ref{thm:essential}.

\section{From a functorial atlas to gluing data}
Assume temporarily that a finite affine atlas of $F$ is given:
\[
 U_1,\ldots,U_n:F\to\Zar,\qquad h_{A_i}\simeq|U_i|.
\]
Let $x_i\in F(A_i)$ be the resulting generic chart point, and put
\[
 d_{ij}=U_j(x_i)\in\Zar(A_i),\qquad X_i=\Spec A_i.
\]
The preceding pullback calculation identifies $|U_i\wedge U_j|$ with
the functor of points of $X_i|d_{ij}$. It also identifies it with that
of $X_j|d_{ji}$. These open subspaces are qcqs, since compact opens of
qcqs schemes have finite affine covers. Full faithfulness therefore
produces transition isomorphisms
\[
 \theta_{ij}:X_i|d_{ij}\simeq X_j|d_{ji}.
\]
First the classifiers of a third compact open show that $\theta_{ij}$ takes
$X_i|(d_{ij}\wedge d_{ik})$ to
$X_j|(d_{ji}\wedge d_{jk})$. On every affine probe the two inverse
images of $U_k$ agree, since the transition retains the same point of
$F$. Affine probes detect equality of compact opens of qcqs spaces:
on an affine atlas one tests with its inverse canonical affine
isomorphisms, then uses lattice separation on the atlas. This proves
the domain equation.

Now restrict the transitions to those identified triple domains. The
two composites into the third chart give the same point of $F$ on
every affine probe. A chart inclusion into $F$ is an embedding, so the
composites agree as chart points; faithfulness makes them equal as
geometric maps. This is the cocycle law on its specified domain.
Identity and inverse laws follow by the same argument. No choice of
generators for the overlap opens is involved.

\subsection*{The restricted cocycle and its section equation}
Here are the domains and restrictions in the preceding argument explicitly.
Use the contravariant indices of the gluing construction: put
$Y_{ij}=X_j|d_{ji}$, let $h_{ij}:Y_{ij}\to X_i$ be the transition
followed by inclusion, and write $t_{ij}(a)$ for its inverse image,
viewed in $L_j$. Its image lies below $d_{ji}$. The transition square
says that $h_{ij}$ and the inclusion of $Y_{ij}$ into $X_j$ induce the
same point of $F$ on every affine probe.

For three indices set $q=t_{jk}(d_{ji})$. The already proved domain
equation, after including the slice lattice into $L_k$, is
\[
 q=d_{kj}\wedge d_{ki}.
\]
In particular $q\leq d_{kj},d_{ki}$. Restrict $h_{jk}$ to $X_k|q$;
its inverse image of $d_{ji}$ is top, so it factors uniquely through
$Y_{ij}$. Call this factor $a$. There are now two actual maps with the
same source and target,
\[
 h_{ij}a,\quad h_{ik}\bigl(X_k|q\longrightarrow Y_{ik}\bigr)
       :X_k|q\longrightarrow X_i.
\]
Their points of $F$ agree by applying the transition square twice on
the left and once on the right; the factorisation triangles and
associativity identify the intermediate chart point. The final chart
inclusion is injective on affine points. Since $X_k|q$ is qcqs, affine
probe separation makes the displayed maps equal.

Taking inverse images of $b\in L_i$ yields
\[
 t_{jk}(t_{ij}(b))\wedge q=t_{ik}(b)\wedge q.
\]
The left meet simplifies by monotonicity from $t_{ij}(b)\leq d_{ji}$.
On the right use $q=d_{kj}\wedge d_{ki}$ and
$t_{ik}(b)\leq d_{ki}$. This proves the exact lattice cocycle
$t_{jk}(t_{ij}(b))=d_{kj}\wedge t_{ik}(b)$.

For sections, apply equality of the two maps after restricting their
images to the common open $w=t_{jk}(t_{ij}(b))\leq q$. On the left,
naturality of $h_{jk}$ and composition of restrictions remove the
restriction introduced by the factorisation. On the right, composition
of restrictions gives the restriction of $c_{ik,b}(s)$ directly to $w$.
Thus
\[
 c_{jk,t_{ij}(b)}(c_{ij,b}(s))
       =c_{ik,b}(s)|_w.
\]
In the overlap slice $Y_{jk}$, the restriction on the right is the
restriction first to $t_{ik}(b)\wedge d_{kj}$ and then along its equality
with $w$. Composition of restrictions identifies these expressions.
This is precisely compatibility in the equaliser defining the glued
section ring; no equality of section carriers is presumed.

For a repeated chart, $d_{ii}=\top$: its generic point belongs to $U_i$.
The transition square and chart injectivity give $h_{ii}$ equal to the
open inclusion. Inverse images and restriction of section images give
$t_{ii}(b)=b$ and the section identity. Applying the same point argument
to the reverse transition gives its composite with the forward
transition equal to identity; cancellation of the final open inclusion
gives the inverse equation in the overlap space itself.

\paragraph{Review before formalisation.}
The source $X_k|q$ was fixed before comparing either composite.
Factoring through $Y_{ij}$ uses $q=t_{jk}(d_{ji})$, whereas restricting
the direct transition uses $q\leq d_{ki}$; these are distinct uses of
the domain equation. Section equality is always stated after restriction
to a common open, so no transport of ring carriers is implicit. The
triple source is qcqs by restriction of the affine $X_k$, even when it
is not affine. Finally an actual atlas still includes only merely given
affine witnesses: finite choice at universe $\ell+1$ supplies all these
witnesses merely, and both this truncation and the atlas truncation are
eliminated only into mere representability. Empty finite families are
allowed throughout.

\section{A small finite gluing construction}
Here is the finite gluing theorem needed in the reconstruction, in the
ringed-lattice approach of \cite{CLS}. Let
$X_i$ have small lattices $L_i$, let $d_{ii}=\top$, and suppose the
transition isomorphisms above are given with the domain and cocycle
laws. Define
\[
 L=\left\{(v_i)\in\prod_i L_i\ \middle|\
   \theta_{ij}^*(v_j\wedge d_{ji})=v_i\wedge d_{ij}
   \text{ for all }i,j\right\}.
\]
Both sides are viewed in $L_i$. Operations are componentwise. The
transition maps preserve the operations on the restricted lattices,
so the displayed equations are stable under joins, meets, bottom and
top. The equations are propositions and $L$ is a small set.

Here is a precise form of the transition laws used in this construction.
Restriction followed by transport defines a bounded lattice map
\[
 r_{ij}:L_i\longrightarrow L_j|d_{ji}.
\]
Write $t_{ij}$ for its underlying function into $L_j$. The identity and
domain-sensitive cocycle equations are
\[
 t_{ii}(a)=a,\qquad
 t_{jk}(t_{ij}(a))=d_{kj}\wedge t_{ik}(a).
\]
In particular $t_{ji}(t_{ij}(a))=d_{ij}\wedge a$.
These equations include the domains of the compositions; an equation of
transitions on an unspecified triple intersection would not suffice.
The corresponding local ring maps are
\[
 c_{ij,a}:\mathcal O_i(a)\longrightarrow
               \mathcal O_j(t_{ij}(a)).
\]
Their identity equation says that $c_{ii,a}(s)$ restricts back to $s$
along $t_{ii}(a)=a$. Their cocycle equation says
\[
 c_{jk,t_{ij}(a)}(c_{ij,a}(s))
   =c_{ik,a}(s)|_{\,t_{jk}(t_{ij}(a))},
\]
where the restriction is defined by the preceding domain equation.
Thus the transported sections form a compatible tuple. These lattice
and section equations are sufficient for the chart argument below.

For $v=(v_i)$, define $\mathcal O(v)$ to be the ring of tuples
$s_i\in\mathcal O_i(v_i)$ whose restrictions agree under the transition
maps on $v_i\wedge d_{ij}$. This is an equaliser of maps of products
of small rings. Restriction and ring operations are componentwise.
Sheaf gluing is componentwise; agreement of the glued components follows
by separation on the overlap. Invertibility opens are
\[
 D_v(s)=(D_{v_i}(s_i))_i.
\]
Locality of the transitions makes this an element of $L$. A compatible
tuple is a unit precisely when its components are units: their inverses
are compatible by uniqueness of inverses. This proves the defining
property of the invertibility open, and the support laws hold
componentwise.

The $i$th chart open has $j$th component $d_{ji}$. Its membership in $L$
uses the domain equations of the transitions, and the charts cover
because $d_{jj}=\top$. Projection identifies its lower interval with
$L_i$. The inverse takes $a\in L_i$ to the family whose $j$th component
is the image of $a\wedge d_{ij}$ under the inverse transition on opens.
The domain and cocycle equations prove compatibility of that family.
They also show that a section tuple on this chart is uniquely determined
by its $i$th component. We obtain $X|e_i\simeq X_i$, with the prescribed
transition maps. Thus $X$ is a small qcqs scheme.

This construction uses only finite products and equaliser conditions in
$\UU_\ell$. In particular it does not construct a large locally ringed
lattice and then silently resize it. For an empty atlas the empty
products give the one-element lattice and zero section ring, hence the
empty scheme; its value on the zero test ring is a singleton.

\section{Identifying the represented functor}
We need the following descent principle in the category of Zariski
sheaves. If $V_i$ is a finite compact-open cover of a functor $G$, and
$H$ is a Zariski sheaf, then compatible natural maps $|V_i|\to H$ glue
uniquely to a map $G\to H$.

To define its value at $x\in G(B)$, choose finite basic presentations
of the compact opens $V_i(x)$, whose union covers $\Spec B$. On each
basic piece, $x$ becomes an element of the relevant $|V_i|$.
Apply the given local map, and use compatibility on intersections to
glue in $H$. Uniqueness is sheaf separation. The type of such an amalgam
is a proposition before a merely existing presentation is eliminated.
Naturality follows by restricting this characterisation along a ring
map and testing on the pulled-back basic cover. This proves the descent
principle without global choices.

More explicitly, a value $y\in H(B)$ fits $x\in G(B)$ when, for every
$h:B\to C$ and every $i$ with $G(h)(x)\in|V_i|(C)$,
\[
 H(h)(y)=\alpha_i(G(h)(x)).
\]
Two fitting values agree on any finite basic refinement of the $V_i(x)$,
hence agree by separation. Existence is first proved using an actual
refinement and then descended to this proposition of fitting candidates.
The base-change square
\[
\begin{array}{ccc}
 B&\longrightarrow&B[1/a]\\
 \downarrow&&\downarrow\\
 C&\longrightarrow&C[1/h(a)]
\end{array}
\]
comes directly from the universal property of localisation. It compares
the local amalgam equations after changing the test ring. This proves
fitting for all $h$, and fitting in turn proves naturality by uniqueness.
Consequently the type of natural extensions is contractible.

Equality of maps into a sheaf can therefore be tested on a finite
compact-open cover. In particular compatible local natural equivalences,
whose inverses are also compatible on overlaps, glue to a natural
equivalence: the two composites are identities on each member of the
respective covers, hence identities globally.

The charts of the constructed $X$ and those of $F$ now give maps between
$F$ and $h_X$: apply the descent principle in each direction. The overlap
conditions are exactly the transition equations already established.
The two composites restrict to identities on all charts, hence are
identities by uniqueness. Therefore $F\simeq h_X$.

Finally, the atlas of $F$ merely exists. An actual atlas gives a witness
of representation; insert it into
\[
 \tr{\sum_X\bigl(\text{$X$ is small and qcqs}\bigr)\times(F\simeq h_X)}.
\]
This is a proposition, so the mere atlas can be eliminated into it.
We do not claim an untruncated canonical geometric object without the
additional univalence theorem for the geometric category.

\section{Over a base ring}
Fix a small commutative ring $k$ with set carrier. Relative functors are
native functors on actual $k$-algebra tests $(B,\iota:k\to B)$ and
algebra maps. Their sheaf condition uses finite basic algebra diagrams;
their compact opens are natural maps to the restriction of
$\operatorname{Zar}$, and their affine charts are represented by small
$k$-algebras. These are intrinsic relative conditions, not definitions
in terms of an absolute total functor.

Define
\[
 (\mathsf T F)(B)=\sum_{\iota:\Hom(k,B)}F(B,\iota).
\]
For $h:B\to C$ let $c_{h,\iota}:(B,\iota)\to(C,h\iota)$ be its
canonical algebra lift. The total action sends $(\iota,x)$ to
$(h\iota,F(c_{h,\iota})(x))$. The dependent comparison needed for every
subsequent argument is
\[
 \operatorname{transport}_{j\mapsto F(C,j)}
       (p,F(c_{h,\iota})(x))=F(h,p)(x),
 \qquad p:h\iota=\kappa.
\]
Path induction proves it: at reflexivity the algebra maps have the same
underlying ring map, and their commuting proofs are propositions.
Consequently total reindexing along any actual algebra map $a$ satisfies
\[
 (\mathsf T F)(\underline a)(\iota,x)=(\kappa,F(a)(x)).
\]
The total identity and composition laws follow from this whole-pair
equation and the native functor laws. Projection gives a natural map
$\pi_F:\mathsf T F\to h_k$.

\begin{theorem}[Intrinsic relative totalisation]
A native relative functor $F$ is a Zariski sheaf if and only if
$\mathsf T F$ is an absolute Zariski sheaf. It is functorially qcqs if
and only if $\mathsf T F$ is absolutely functorially qcqs.
\end{theorem}
\begin{proof}
For a compatible total cover family $(\iota_i,x_i)$, first glue the base
maps in the representable sheaf $h_k$. Write the resulting equations as
$p_i:\lambda_i\iota=\iota_i$. Transport $x_i$ backwards along $p_i$
into $F(B_i,\lambda_i\iota)$. Their pairs equal the original total
points. The whole-pair reindexing equation above transfers compatibility;
injectivity of pairing at the fixed overlap base gives native relative
compatibility. Glue the rebased values in $F$. For uniqueness, any other
total amalgam has the same base map by uniqueness in $h_k$; rebase its
value along that equation and apply relative uniqueness.
Conversely, total gluing of paired native values and separation in $h_k$
rebase the glued point into the specified fibre. Total uniqueness and
fixed-base injectivity give relative uniqueness. Empty covers are included.

Compact opens and their cover equations pass between the relative and
total functors by evaluation at $(\iota,x)$. Reordering the membership
proof gives $\mathsf T|U|\simeq|\mathsf T U|$. For an algebra
$(A,\alpha)$, forgetting the structure map gives
$\mathsf T h^k_{(A,\alpha)}\simeq h_A$: that structure is forced to be
$h\alpha$ by its underlying ring map $h:A\to B$. Thus relative affine
atlases totalise. Conversely a chart $e:h_A\simeq\mathsf T G$ supplies
$\alpha:k\to A$ from the base component of $e_A(1_A)$. Naturality says
that the base of $e_B(h)$ is $h\alpha$. Transporting along that equation
restricts $e$ to a native relative affine equivalence. Applying this to
the total chart realisations recovers a relative affine atlas. Atlas
and affine-witness truncations are eliminated only into propositions.
\end{proof}

\begin{theorem}[Relative representation and categorical comparison]
A native functor $F$ on small $k$-algebras is functorially qcqs if and
only if there merely exist a small geometric qcqs space $X$, a map
$s:X\to\Spec k$, and a native natural equivalence
\[
 F(B,\iota)\simeq\Hom_{\Spec k}(\Spec(B,\iota),(X,s)).
\]
The relative functor of points is an equivalence of univalent categories
between geometric qcqs schemes over $\Spec k$ and intrinsic relative
functorial qcqs schemes. It has an explicit inverse functor and both
natural isomorphisms.
\end{theorem}
\begin{proof}
Totalise $F$ and apply the absolute representation theorem. For an actual
representation $e:\mathsf T F\simeq h_X$, compose $e^{-1}$ with $\pi_F$
to get $q:h_X\to h_k$. The affine comparison and absolute full
faithfulness give $s:X\to\Spec k$ and the equation
$s p=\Spec(q_B(p))$ on every affine probe. Taking fibres over $\iota$
therefore gives geometric points over $k$. On the original side the
fibre is
\[
 \sum_{\kappa:k\to B}\sum_{x:F(B,\kappa)}(\kappa=\iota)
       \simeq F(B,\iota),
\]
where the equivalence transports $x$ along the displayed path. Path
induction gives its two inverse laws. Restricting $e$ to fibres and
composing these equivalences gives the native relative representation.
The mere absolute representation is eliminated only into mere relative
representability. Conversely, the total relative geometric functor is
$h_X$, by the canonical lift and uniqueness of the base map of an
ordinary affine point. The converse totalisation criterion and invariance
under relative natural equivalence prove the necessary direction.

The geometric category over $k$ is a univalent slice. Native relative
functors form a univalent category by the value-family and action
identity argument; the relative qcqs condition is a proposition.
Existing relative full faithfulness and the just-proved relative
essential image give the categorical comparison by the same
propositional-fibre argument as in the next section.
\end{proof}

The companion note \emph{Relative totalisation over a base ring: a proof
with explicit changes of structure map} gives the detailed argument,
its review, and the declaration-by-declaration interfaces used to handle
the dependent types. It was written before formalising this continuation.

\section{Equivalence of univalent categories}
This section gives the categorical conclusion, including its univalence
arguments. The structured identity principles, category packaging,
comparison functor and explicit quasi-inverse are now checked, as well
as the preceding absolute Hom--Nat and essential-image theorems.
The informal proof and its review were written before formalisation.

Let $\mathcal C$ have objects small locally ringed lattices $X$ together
with a proof that $X$ is qcqs, and local geometric maps as arrows. Let
$\mathcal D$ have objects small-set-valued functors $F$ together with a
proof that $F$ is functorially qcqs, and natural transformations as arrows.
Their identities and compositions are the existing geometric and natural
ones. Geometric maps and natural transformations form sets. Both object
types lie in $\UU_{\ell+1}$; geometric homs lie in $\UU_\ell$, whereas
natural transformations are initially in $\UU_{\ell+1}$. All categories
can consequently be considered at the common bound $\ell+1$.
No resizing of the natural-transformation types is needed.

The functor of points is
\[
 H:\mathcal C\longrightarrow\mathcal D,\qquad
 H(X)=h_X,\quad H(f)_B(p)=f\circ p.
\]
The geometric atlas theorem supplies its object property. The unit and
associativity equations prove that it is a functor. The checked
Hom--Nat theorem says that $H$ is fully faithful, and the checked
essential-image theorem says that it is merely essentially surjective.

\begin{lemma}[Univalence of the geometric category]
For $X,Y\in\mathcal C$, the canonical map
\[
 (X=Y)\longrightarrow\operatorname{Iso}_{\mathcal C}(X,Y)
\]
is an equivalence.
\end{lemma}
\begin{proof}
First consider commutative rings with set carriers and bounded
distributive lattices. Univalence identifies equality of their carriers
with equivalence of carriers. Under this identification, transport of
operations is exactly preservation of those operations. Function
extensionality identifies the corresponding equations of operations;
the equational axioms and sethood proofs are propositions. This proves
the structured identity principle for these two algebraic structures:
paths correspond to algebra isomorphisms, with the correspondence being
the canonical one at reflexivity.

For a fixed lattice $L$, consider presheaves of rings with set carriers.
An identity of their ring families corresponds to a family of ring
isomorphisms. After transporting along that identity, equality of the
restriction maps is exactly naturality of the family of isomorphisms.
Restriction laws and the sheaf condition are propositions. Thus equality
of sheaves of rings on $L$ corresponds to natural ring isomorphisms.
This assertion follows by equivalence induction on the ring families:
in the identity case naturality states equality of restrictions, and
function extensionality gives the equality of restriction families.

Now an isomorphism $f:X\to Y$ gives a lattice isomorphism
$f^*:L_Y\to L_X$ and ring isomorphisms
\[
 f^\sharp_u:\mathcal O_Y(u)\longrightarrow
                  \mathcal O_X(f^*u)
\]
compatible with restrictions. Conversely, such data give an isomorphism
of ringed lattices. Any such isomorphism between locally ringed lattices
is local: ring isomorphisms preserve and reflect units, so the defining
universal property of invertibility opens identifies the two opens.
Its inverse is local by the same argument.

Use the inverse lattice isomorphism to identify $L_X$ with $L_Y$.
With this base fixed, the preceding sheaf identity principle identifies
the sheaves using $f^\sharp$ (in the appropriate inverse orientation).
For a fixed sheaf, the type of invertibility maps is a proposition by
uniqueness of invertibility opens; its support laws are propositions too.
Consequently the locally ringed structures are identified, including
these fields. Finally the qcqs witnesses lie in a proposition, so they
add no further identity data.

These identifications are equivalences at each stage of the dependent
record, rather than just a construction of some path from an isomorphism.
They agree with the identity isomorphism at reflexivity. Path induction
therefore identifies their composite with the canonical map from paths
to geometric isomorphisms. This proves the required univalence.
\end{proof}

\begin{lemma}[Univalence of the functorial category]
For $F,G\in\mathcal D$, the canonical map
\[
 (F=G)\longrightarrow\operatorname{Iso}_{\mathcal D}(F,G)
\]
is an equivalence.
\end{lemma}
\begin{proof}
A categorical isomorphism of set-valued functors is exactly a natural
transformation with equivalence components: its pointwise inverse is
natural by injectivity of the forward components. Apply univalence
pointwise to identify the value sets along those component equivalences,
and function extensionality to identify the value families. After this
identification, naturality is exactly equality of the reindexing maps.
The sethood proofs and functor laws are propositions. This gives the
structured identity principle for set-valued functors, with canonical
map the identity-to-natural-isomorphism map. The functorial qcqs condition
is a proposition, so restriction to that full subcategory preserves
univalence.
\end{proof}

\begin{theorem}[Categorical comparison]
The functor $H:\mathcal C\to\mathcal D$ is an equivalence of univalent
categories. Its map on object types is an equivalence.
\end{theorem}
\begin{proof}
Full faithfulness induces equivalences on isomorphism types. Indeed,
lift an isomorphism and its inverse by fullness; their inverse equations
hold by faithfulness. These lifts are unique by the equivalences on hom
sets. For $X,Y\in\mathcal C$ we therefore have the commutative square
\[
\begin{array}{ccc}
 (X=Y)&\xrightarrow{\operatorname{ap}H}&(HX=HY)\\
 \downarrow\scriptstyle{\simeq}&&\downarrow\scriptstyle{\simeq}\\
 \operatorname{Iso}_{\mathcal C}(X,Y)&\xrightarrow{\simeq}&
 \operatorname{Iso}_{\mathcal D}(HX,HY).
\end{array}
\]
Commutativity is path induction; the vertical maps are univalence and
the bottom map is the preceding equivalence on isomorphisms. Thus
$\operatorname{ap}H$ is an equivalence: the object map of $H$ is an
embedding.

For $F\in\mathcal D$, its fibre
\[
 P_F=\sum_{X:\mathcal C}(HX=F)
\]
is consequently a proposition. The essential-image theorem gives merely
an $X$ and a natural equivalence $F\simeq HX$. Invert that equivalence
and apply univalence of $\mathcal D$ to get $\tr{P_F}$. Since $P_F$ is a
proposition, eliminate this truncation into $P_F$ itself. It is therefore
inhabited and contractible. Every fibre of the object map is contractible,
so $H$ is an equivalence on objects. This is the precise step that removes
the truncation; it uses univalence and full faithfulness, not a choice
of affine atlases.

For completeness, its categorical inverse can be read off from these
fibres. Write their centres as $(KF,\epsilon_F:HKF=F)$ and let
$e_F:HKF\cong F$ be the resulting isomorphism. For
$\alpha:F\to G$, define $K\alpha$ as the unique geometric arrow with
\[
 H(K\alpha)=e_G^{-1}\circ\alpha\circ e_F.
\]
Existence and uniqueness use the hom equivalence. Faithfulness proves
the identity and composition laws. The $e_F$ form a natural isomorphism
$HK\cong1_{\mathcal D}$ by this defining equation. Lifting $e_{HX}$
gives a natural isomorphism $KH\cong1_{\mathcal C}$; it is an
isomorphism since full faithfulness reflects isomorphisms. Hence $K$ is
a quasi-inverse functor, as claimed.
\end{proof}

\paragraph{Review.}
The geometric univalence proof includes the ring-family and restriction
maps, not just the lattice. The inverse orientation of inverse images
is respected. Invertibility maps are unique by their universal property;
they must not be treated as arbitrary extra operations. Adding a merely
existing affine atlas as a property introduces no chosen cover.
Both univalence statements assert equivalence of identity types with
isomorphisms, rather than only that isomorphic objects can be identified.
In particular automorphisms remain as loops of the object types.
The representation fibres are propositions only after full faithfulness
and univalence have been established. The resulting removal of
truncation uses no axiom of choice and entails no universe resizing.

\section{Review and connection with the implementation}
The proof separates four interfaces: functors on all small tests and
Yoneda; compact opens and their realisations; finite compact-cover
descent for maps into sheaves; and finite geometric gluing with cocycles.
Compact-open atlases account for quasi-separatedness. The finite tuple
construction accounts for smallness. Mere existence of the atlas is
used only after truncating the representation statement. None of these
points is supplied by full faithfulness alone.

\begingroup\raggedright
The existing checked development supplies full faithfulness, geometric
finite-cover descent, restriction to compact opens, and their qcqs
property. The first new checkpoint now supplies the arbitrary-functor
and compact-open interfaces described above. In particular,
\texttt{TestFunctors}, \texttt{TestYoneda}, and \texttt{TestEquivalences}
give the functor interface on all small rings; \texttt{ZariskiFunctor}
gives the classifier and its comparison with inverse images under
$\Spec$. \texttt{CompactOpens}, \texttt{OpenRealisation}, and
\texttt{CompactPullback} give the lattice, realisation and affine
pullback formula. \texttt{GeometricOpens} proves the realisation equation
for geometric opens.

\texttt{TestSheaves} states the finite basic sheaf condition, including
empty covers. \texttt{RepresentableSheaf} proves it for all affine
representables using the earlier localisation gluing theorem.
\texttt{ZariskiSeparated} proves separation of the classifier;
\texttt{OpenSheaf} proves that compact-open realisations of sheaves are
sheaves. \texttt{WholeRealisation} and \texttt{SheafTransfer} supply the
whole-open equivalence and invariance under natural equivalence.
\texttt{FunctorialQcqs} gives the intrinsic definition.
\texttt{BasicAffineFunctor} identifies $|D(a)|$ with $h_{A[1/a]}$;
\texttt{CompactAffineCover} proves that every compact open of an affine
representable is a functorial qcqs scheme.

\texttt{GeometricEquivalence} and \texttt{AffineFunctor} transfer
geometric isomorphisms and affine objects to functors.
\texttt{SpaceGluing} states geometric finite-cover descent directly in
terms of maps agreeing on intersections. \texttt{OpenMonic} and
\texttt{BasicGeometry} identify the geometric basic-open restriction
diagram with the localisation diagram. \texttt{GeometricSheaf} then
proves sheafhood of $h_X$ for every locally ringed lattice $X$.
Together with the atlas construction, \texttt{GeometricAtlas.qcqs}
proves that every geometric qcqs scheme gives a functorial qcqs scheme.

\texttt{FiniteLatticeGlue} and \texttt{FiniteRingGlue} construct the
small compatible-tuple lattice and rings with their universal
properties; units are detected coordinatewise.
\texttt{FiniteRingedGlue} supplies the restriction maps and sheaf proof,
and \texttt{FiniteLocalGlue} supplies the locally ringed structure,
local projection maps, and their equality on overlaps.
\texttt{FiniteTransitionLattice} proves the covering equation and the
lattice chart theorem under the displayed domain-sensitive laws.
\texttt{FiniteTransitionGlue} proves the section-ring chart theorem,
the geometric chart isomorphisms, and the qcqs conclusion for affine
input charts. The extraction of these hypotheses from a functorial atlas
is now checked in the atlas modules described below.

\texttt{LocalisationTests} and \texttt{CompactMembership} give the
base-change square and local membership criterion used in descent.
\texttt{CompactCoverDescent} proves contractibility of natural extensions,
with no sheaf assumption on the source. \texttt{CompactCoverSeparation}
gives local detection of equality, and \texttt{LocalFunctorIso} glues
compatible local natural equivalences of sheaves.

The overlap construction is also checked. \texttt{RealisationTransfer}
restricts a natural equivalence using a membership equivalence;
\texttt{OpenIntersections} and \texttt{OpenSymmetry} identify realised
intersections while retaining their underlying point.
\texttt{AffineOpenRealisation} represents every compact open of $h_A$
by the corresponding open subspace of $\Spec A$.
\texttt{AtlasOverlap} therefore represents affine-chart overlaps by
qcqs spaces and records their inclusion into the original functor.
\texttt{GeometricRecovery} recovers geometric isomorphisms from natural
equivalences, and \texttt{AtlasTransition} gives the transition
isomorphisms and their equations over that functor.
\texttt{GeometricOpenSeparation} shows that affine probes detect
equality of compact opens of qcqs spaces. Using it,
\texttt{AtlasDomains} proves that a transition identifies the two
inverse images of any third compact open, hence the triple-overlap
domain equation.

One change of presentation is treated explicitly. \texttt{SpacePaths}
transports a diagram along equalities of its spaces and proves that
composition and isomorphisms are preserved. The basic-open diagram is
checked at its defining presentation and moved as a whole. This avoids
rechecking the component proofs of the lattice and sheaf when packaging
the diagram in geometric notation.

The absolute essential-image theorem is now checked.
\texttt{AtlasInverse} proves the diagonal and reverse-transition laws.
\texttt{OpenFactor} restricts and factors maps through an open with an
explicit lattice and section interface. \texttt{AtlasData} packages the
actual finite chart data and the domain equations.
\texttt{AtlasCocycle} compares the two actual maps on $X_k|q$;
\texttt{AtlasLaws} and \texttt{AtlasSections} extract exactly the lattice
and section hypotheses of finite gluing. \texttt{AtlasReconstruction}
applies that theorem and records the new chart classifiers.
\texttt{AtlasCompatibility} proves compatibility in both directions,
and \texttt{AtlasRealisation} applies the local-to-global equivalence
result.

For clarity, the compatibility argument has two directions. If points
$p$ in chart $i$ and $r$ in chart $j$ have the same $F$-point, the
classifier puts $r$ in $d_{ji}$. Factor $r$ through that overlap. Its
transition image equals $p$ by injectivity of the original chart.
The equality of the two overlap maps into the glued space then gives
equal geometric points. Conversely, if their geometric points agree,
the reconstructed classifier puts $r$ in $d_{ji}$. The same overlap
square and injectivity of the reconstructed chart identify its
transition image with $p$; the transition trace gives equal $F$-points.
Thus neither inverse compatibility nor agreement of classifiers is
assumed in applying the local-to-global equivalence.

\texttt{FiniteLargeChoice} proves finite choice for large affine witnesses
indexed by the existing small finite type; finite types at different
universe levels are not silently identified. \texttt{FunctorialTransfer}
transfers sheafhood and affine atlases across a natural equivalence.
\texttt{EssentialImage.reconstruct} eliminates the finite witnesses and
the atlas into mere representability; \texttt{recognise} proves the
converse. \texttt{equivalence} identifies the two propositions in
Theorem~\ref{thm:essential}. The relative totalisation and representation theorems of Section~7 are
now checked too, including the equivalence of relative univalent categories.
The absolute categorical comparison of Section~8 remains the input to
that reduction.

The informal restricted-cocycle proof and review were written before
formalisation. During the opening audit the section-composite argument
was separated into its stated intermediate section images, and the
factor map was made opaque after its lattice and restriction equations.
The reconstruction triangle is checked at its defining glued space.
These changes keep the formal proof aligned with its mathematical steps.
The categorical continuation reuses \texttt{SipGeneric} and the existing
\texttt{Category}, \texttt{Arrows}, \texttt{Functor}, and
\texttt{EquivalenceLemmas} interfaces without changing them.
\texttt{RingIdentity} and \texttt{LatticeIdentity} apply the generic
structure identity theorem to the actual algebra records.
\texttt{FamilyIdentity} and \texttt{DecoratedIdentity} pass identity
systems to families, structure, and propositional properties.
\texttt{PresheafIdentity} and \texttt{LocalIdentity} apply these results
to rings with restriction maps and their locally ringed properties.
\texttt{SpaceIdentity} handles the dependent lattice/structure pair;
\texttt{SpaceIsoIdentity} identifies the structural presentation of
an actual ringed-lattice isomorphism with its original ringed map.
\texttt{GeometricCategory} uses this bridge and the categorical
singleton criterion to prove univalence for all spaces.

\texttt{SetIdentity} and \texttt{FunctorIdentity} prove the analogous
identity principle for the existing set-valued functors;
\texttt{FunctorCategory} proves categorical univalence.
\texttt{FullSubcategory} restricts both categories to their qcqs
properties. \texttt{WeakCategoryEquivalence} proves the object
embedding and removes the truncation of essential surjectivity using
\texttt{EmbeddingFibres}. \texttt{CategoryQuasiInverse} constructs the
inverse functor and its two natural isomorphisms, with checked identity,
composition, and naturality laws.

\texttt{QcqsCategories.equivalence} assembles the categorical equivalence
for the existing functor of points. Its source and target are actual
univalent category records; their univalence proofs are constructed in
this continuation, rather than parameters assumed by the comparison.
\texttt{object-equivalence} separately records the equivalence on object
types. These conclusions use the existing universe bounds and no axiom
of choice or resizing.

The relative continuation follows the detailed companion note.
\texttt{RelativeTotal.change-target} states the dependent transport
comparison, and \texttt{on-algebra} compares the whole total point for
an arbitrary algebra map. \texttt{RelativeCoverFibre} separates pair
equalities from native fibre equalities; \texttt{RelativeBaseGlue},
\texttt{RelativeTotalSheaf}, and \texttt{RelativeSheafOfTotal} implement
the two sheaf arguments. \texttt{RelativeTotalisation} exposes both the
sheaf and qcqs equivalences, with intrinsic native definitions.
\texttt{RelativeEssentialImage} proves relative reconstruction and its
converse. \texttt{RelativeCategories.Over[K].equivalence} packages the
relative categorical comparison, using \texttt{SliceCategory} and
\texttt{RelativeFunctorCategory} for univalence.

The complete qcqs tree now checks with 249 modules and 9,250 declarations,
zero failures. All 3,752 lists in the combined absolute, category, and
relative opening audit report no extra openings. No construction proof
of function extensionality, localisation, or quotients is reopened by
the relative clients. The detailed note names the intermediate base
paths, rebased values, and pointwise squares that keep the checker
aligned with the mathematical argument.

The approach is the compact-open formulation of
\cite[Sections 4--5]{Zeuner} and \cite[Sections 5--6]{ZH}; its finite
small gluing construction replaces the universe-lifted realisation route.
\par\endgroup

\begin{thebibliography}{9}
\bibitem{CLS} Thierry Coquand, Henri Lombardi and Peter Schuster,
\emph{Spectral schemes as ringed lattices}, Annals of Mathematics and
Artificial Intelligence \textbf{56} (2009), 339--360.
\bibitem{Zeuner} Max Zeuner, \emph{Univalent Foundations of Constructive
Algebraic Geometry}, Sections 4--5. Local source: \texttt{SCHEME/zeuner1.pdf}.
\bibitem{ZH} Max Zeuner and Matthias Hutzler, \emph{The Functor of Points
Approach to Schemes in Cubical Agda}, Sections 5--6.
Local source: \texttt{SCHEME/zeuner3.pdf}.
\end{thebibliography}
\end{document}
